\section{Step 3 Stop before and after a call}\label{transpose-s03}

The supplied caller chooses a function address in \texttt{\$t9}, then calls through that address at \texttt{transpose\_begin}. It places the two full pointers in the first two argument registers and the dimensions in the next two. The continuation at \texttt{transpose\_done} records the count only after the function returns. This separates the function measurement from the caller's store.

\subsection{Write or inspect}
Set source breakpoints at \texttt{transpose\_begin} and \texttt{transpose\_done}. Inspect \texttt{a0} through \texttt{a3} at entry. The stubs return immediately; do not interpret their output as a completed transpose.

\subsection{Expected observation}
The default entry arguments name the two matrix bases and dimensions three and five. The first breakpoint precedes the call, and the second precedes the result store.

\Needspace{11\baselineskip}
\subsection{Check your understanding}
When should a function-only cache measurement stop?

\begin{itemize}
  \item \textbf{A} After the caller stores the returned count.
  \item \textbf{B} After the exit syscall.
  \item \textbf{C} Before the continuation executes its result store.
\end{itemize}
